\documentclass[10pt,a4paper]{amsart}
\usepackage[margin=19mm]{geometry}
\usepackage{amsmath,amssymb,mathtools}
\usepackage[scr=rsfs,cal=euler]{mathalfa}
\usepackage{xfrac}
\usepackage[unicode,hidelinks,bookmarks=false]{hyperref}
\newcommand{\ie}{i.e.\ }
\newcommand{\cf}{cf.\ }
\newcommand{\resp}{respectively\ }
\newcommand{\Subsection}{\subsection}
\providecommand{\nobreakdash}{\nobreak\hbox{-}}
\setlength{\emergencystretch}{2em}
\multlinegap0pt
\usepackage{bm}
\usepackage{bbm}
\usepackage{dsfont}
\usepackage{xspace}
\usepackage{cancel}
\usepackage{mathtools}
\usepackage[shortlabels]{enumitem}
\usepackage{amsthm}
\makeatletter
\@ifundefined{proposition}{\newtheorem{proposition}{Proposition}}{}
\makeatother
\title[Computations of $\tilde{A}\_2$ Bruhat intervals via shadows]{Computations of $\tilde{A}\_2$ Bruhat intervals via shadows}
\author{Megan Masters}
\date{Source: arXiv:2607.21054v1 (CC BY 4.0)}
\begin{document}
\maketitle

\noindent\emph{Source and licence.} Computations of $\tilde{A}_2$ Bruhat intervals via shadows. Version arXiv:2607.21054v1 (CC BY 4.0), distributed under the Creative Commons Attribution 4.0 International licence, as recorded in the arXiv licence metadata for this exact version and shown on the abstract page https://arxiv.org/abs/2607.21054v1. Full rights notice: licenses/arxiv-2607.21054-CC-BY-4.0.txt.\par\smallskip

\noindent\textit{Editorial scope.} Counted unit: the Proposition whose source label is \texttt{regionslemma} in the article (source0.tex lines 684--686, source label \texttt{regionslemma}), delivered together with the article's own complete original proof of it (source0.tex lines 688--800, one \texttt{proof} environment of 113 lines). No mathematics is written, repaired or re-derived: statement and proof are the article's own text line for line. Required context retained uncounted: regions (lines 656--658). Every definition, display and label of those lines is retained unchanged. Omitted from this package and not redistributed: 1--655, 659--683, 687--687, 801--1691. Nothing inside a retained excerpt is omitted or altered: every retained body is delivered exactly as the article's lines are written, so no omission receipt of any kind accompanies this package. Citations: the retained excerpts carry no citation command at all, so no bibliography entry of the article is transcribed and no reference is redistributed. Reference adaptation: none was needed and none was made; every reference inside the delivered unit resolves inside it, so no unresolved reference or citation is produced. Printed slips kept literal: no printed word, symbol, hypothesis or number of the counted statement or of its proof was altered. Layout and class: the article's own class is \texttt{article} with packages including geometry, amsmath, amsthm, bbm, amssymb, graphicx, amsfonts, xcolor, pdfpages, enumerate, hyperref, caption, tikz, booktabs; the wrapper is the lane's pinned \texttt{amsart} editorial wrapper at 10pt on A4, which restates the theorem-like environments the retained text needs and adds the article's own macro definitions that the retained text needs (none), with extra wrapper packages (bm, bbm, dsfont, xspace, cancel, mathtools, enumitem). The delivered numbering is the unit's own. Assets: no retained span contains a figure environment, a graphics-inclusion command, a PDF-page inclusion, a table or a TikZ picture; no image file is copied into this package, the package declares no asset dependency and no image file is redistributed with it. Licence: version arXiv:2607.21054v1 of this article is distributed under the Creative Commons Attribution 4.0 International licence, as recorded in the arXiv licence metadata for this exact version and shown on the abstract page https://arxiv.org/abs/2607.21054v1. A full rights notice is delivered at licenses/arxiv-2607.21054-CC-BY-4.0.txt. Characters: every retained excerpt is pure ASCII, so the delivered engine, XeLaTeX with its default Latin Modern text font, reports no missing character.
\begin{proposition}\label{regions}
    If $w,v\in W$ are two reduced words that differ by a braid relation, then the region produced by the algorithm for $w$ and $v$ is the same region. 
\end{proposition}

\begin{proposition} \label{regionslemma}
    Given two reduced words $w,v\in W$ that differ by a braid relation, we have $n_w=n_v$ and $m_w=m_v$. 
\end{proposition}

\begin{proof}
    The three braid relations we have are 
    \[s_0s_1s_0\longleftrightarrow s_1s_0s_1,\]
    \[s_1s_2s_1\longleftrightarrow s_2s_1s_2,\]
    \[s_2s_0s_2\longleftrightarrow s_0s_2s_0.\]
    We can assume without loss of generality that $w=w_1s_is_{i+1}s_iw_2$ and $v=w_1s_{i+1}s_is_{i+1}w_2$ for some reduced, possibly empty, words $w_1, w_2$ and $i\in \{0,1,2\}$. We proceed by considering cases where $w_1$ and $w_2$ are possibly empty. We apply the algorithm to $w$, outputting coordinates $(n_w,m_w)$. In some of the cases, we further split into two subcases:  the case that the start of the braid relation is contained in an increasing subword and the case that the start is contained in a decreasing subword when the algorithm is applied.
    
    Case 1: $w_1$ and $w_2$ are empty. Then $w=s_is_{i+1}s_i$. So the algorithm tells us that $w$ lies in the $s_i$ region. Now, running the algorithm, we get 
        \[w=\underbrace{s_is_{i+1}}_\uparrow \underbrace{s_i}_\downarrow,\]
        and so $u_w=2$, and $x_w=2$, $y_w=1$. Then applying the braid relation we get 
        \[v=\underbrace{s_{i+1}s_i}_\downarrow \underbrace{s_{i+1}}_\uparrow,\]
        as we know that we are in the $s_i$-region  and we are now starting with $s_{i+1}$, so we start with a decreasing subword. Now we have $u_v=1$, but we now have a word ending in an odd increasing subword. So $x_v=2$, $y_v=1$, and hence $n$ and $m$ are preserved, as, by Proposition \ref{regions}, the region determined by the algorithm is constant under braid relations.

    Case 2: $w_1$ is empty but $w_2$ is nonempty. 
    So there is at least one more element after $s_is_{i+1}s_i$. Then this element must be $s_{i+2}$; if the element was $s_i$ we would immediately see that it is not a reduced word, and if it was $s_{i+1}$ we could apply the braid relation and again see it was not a reduced word. Now we need to consider three subcases, depending on the length of $w_2$ and the element after $s_{i+2}$ in $w_2$.

    \begin{enumerate}[(a)]
        \item $\ell(w_2)=1$, and so $w_2=s_{i+2}$. Then we have
        \[w=\underbrace{s_is_{i+1}}_\uparrow \underbrace{s_is_{i+2}}_\downarrow,\]
        and applying the braid relation we have
        \[v=\underbrace{s_{i+1}s_i}_\downarrow \underbrace{s_{i+1}s_{i+2}}_\uparrow,\]
        as we are now starting with $s_{i+1}$, so we start with a decreasing subword. These both give $x_w=x_v=2$, $y_w=y_v=2$. 
        
        \item $\ell(w_2)>1$, and $w_2=s_{i+2}s_{i+1}\hdots$, with the dots possibly representing an empty string. Then we have
        
        \[w=\underbrace{s_is_{i+1}}_\uparrow \underbrace{s_is_{i+2}s_{i+1}\hdots}_\downarrow\hdots,\]
        and applying the braid relation we have
        \[v=\underbrace{s_{i+1}s_i}_\downarrow \underbrace{s_{i+1}s_{i+2}}_\uparrow\underbrace{s_{i+1}\hdots}_\downarrow\hdots,\]
        as again we are now starting with $s_{i+1}$. Now the number of elements in an increasing subword is preserved, and we have not changed the end subword, and so $n$ and $m$ are preserved.
        \item $\ell(w_2)>1$, and $w_2=s_{i+2}s_i\hdots$, with the dots possibly representing an empty string. Then we have
        
        \[w=\underbrace{s_is_{i+1}}_\uparrow \underbrace{s_is_{i+2}}_\downarrow\underbrace{s_i\hdots}_\uparrow\hdots,\]
        and applying the braid relation we have
        \[v=\underbrace{s_{i+1}s_i}_\downarrow \underbrace{s_{i+1}s_{i+2}s_i\hdots}_\uparrow\hdots,\]
        as again we are now starting with $s_{i+1}$. Now the number of elements in an increasing subword is preserved, and, as we have not changed the parity of the length of a possible end increasing subword, $n$ and $m$ are preserved.
    \end{enumerate}
    Case 3: $w_1$ is nonempty but $w_2$ is empty. As $w_1$ is nonempty,  by the same argument as above, $w_1$ must end in $s_{i+2}$. So our $w$ looks like $w=\hdots s_{i+2}s_is_{i+1}s_i$. We now run the algorithm on this word. We consider two cases: when the $s_{i+2}$ element at the end of the subword $w_1$ is contained in an increasing subword, and when it is contained in a decreasing subword. 
\begin{enumerate}[(i)]
        \item $s_{i+2}$ is contained in an increasing subword. So we have
    \[w=\hdots\underbrace{\hdots s_{i+2}s_is_{i+1}}_\uparrow \underbrace{s_i}_\downarrow.\]
    Let $u$ be the number of elements in increasing subwords. Then we are not ending in an odd increasing subword, so $x_w=u$ and $y_w=\ell(w)-u$.
    
    Applying the braid relation we have 
    \[v=\hdots\underbrace{\hdots s_{i+2}}_\uparrow \underbrace{s_{i+1}s_i}_\downarrow \underbrace{s_{i+1}}_\uparrow.\]
    Now we note that everything up to, and including, our $s_{i+2}$ has not been changed, and so whether it is in an increasing or decreasing subword has not changed. Then, overall, the number of elements in an increasing subword has decreased by one, so we now have $u-1$ elements in increasing subwords. But we are now ending in an odd length increasing subword, so we still have $x_v=u$ and $y_v=\ell(w)-u$. Then our $n$ and $m$ are unchanged. 

    \item $s_{i+2}$ is contained in a decreasing subword. So we have 
    \[w=\hdots\underbrace{\hdots s_{i+2}}_\downarrow \underbrace{s_is_{i+1}}_\uparrow \underbrace{s_i}_\downarrow.\] 
    Applying the braid relation, we get
    \[v=\hdots\underbrace{\hdots s_{i+2}s_{i+1}s_i}_\downarrow \underbrace{s_{i+1}}_\uparrow.\]
    Again, the number of elements in an increasing subword has decreased by one. But our word now ends in an odd increasing subword, and so our $n$ and $m$ are the same before and after the braid relation.
    \end{enumerate}


    Case 4: Both $w_1$ and $w_2$ are nonempty. Then $w_1$ must end with $s_{i+2}$ and $w_2$ must start with $s_{i+2}$. So $w=\hdots s_{i+2}s_is_{i+1}s_is_{i+2}\hdots$, where the dots can be empty. Again, we consider the two cases when the $s_{i+2}$ element at the end of the subword $w_1$ is contained in an increasing subword, and when it is contained in a decreasing subword. Within these cases we need to consider three subcases, just as we did in Case 1, depending on the length of $w_2$. 

    
    \begin{enumerate}[(i)]
        \item Our starting $s_{i+2}$ appears in an increasing subword.
        \begin{enumerate}
        \item $\ell(w_2)=1$, and so $w_2=s_{i+2}$. So we have
        \[w=\hdots\underbrace{\hdots s_{i+2}s_is_{i+1}}_\uparrow\underbrace{s_is_{i+2}}_\downarrow,\]
        which we can apply the braid relation to. Then we get
        \[v=\hdots\underbrace{\hdots s_{i+2}}_\uparrow\underbrace{s_{i+1}s_i}_\downarrow\underbrace{s_{i+1}s_{i+2}}_\uparrow,\]
        and we can clearly see that the number of elements in increasing subwords is preserved. The braid move now means that we end in an increasing subword. But this subword has even length. So overall the values of $n$ and $m$ are preserved under this braid move.

        \item $\ell(w_2)>1$, and $w_2=s_{i+2}s_{i+1}\hdots$, with the dots possibly representing an empty string. Then we have 
        \[w=\hdots\underbrace{\hdots s_{i+2}s_is_{i+1}}_\uparrow\underbrace{s_is_{i+2}s_{i+1}\hdots}_\downarrow\hdots,\]
        and, applying the braid relation, we get
        \[v=\hdots\underbrace{\hdots s_{i+2}}_\uparrow\underbrace{s_{i+1}s_i}_\downarrow\underbrace{s_{i+1}s_{i+2}}_\uparrow\underbrace{s_{i+1}\hdots}_\downarrow\hdots,\]
        which preserves $u$ and whether the word ends in an odd increasing subword. 

        \item $\ell(w_2)>1$, and $w_2=s_{i+2}s_i\hdots$, with the dots possibly representing an empty string. Then we have 
        \[w=\hdots\underbrace{\hdots s_{i+2}s_is_{i+1}}_\uparrow\underbrace{s_is_{i+2}s_{i+1}}_\downarrow\underbrace{s_i\hdots}_\uparrow\hdots,\]
        and, applying the braid relation, we get
        \[v=\hdots\underbrace{\hdots s_{i+2}}_\uparrow\underbrace{s_{i+1}s_i}_\downarrow\underbrace{s_{i+1}s_{i+2}s_i\hdots}_\uparrow\hdots,\]
        which preserves $u$ and the parity of any end subword. 

    
        \end{enumerate}
        \item Our starting $s_{i+2}$ appears in a decreasing subword. Again consider the three subcases.
        \begin{enumerate}
            \item $\ell(w_2)=1$, and so $w_2=s_{i+2}$.
        \[w=\hdots\underbrace{\hdots s_{i+2}}_\downarrow\underbrace{s_is_{i+1}}_\uparrow\underbrace{s_is_{i+2}}_\downarrow,\]
        and applying the braid relation we get
        \[v=\hdots\underbrace{\hdots s_{i+2}s_{i+1}s_i}_\downarrow\underbrace{s_{i+1}s_{i+2}}_\uparrow.\]
        Now clearly the number of elements in an increasing subword is unchanged under the braid move. Also, we are now ending on an increasing subword, but its length is even. So again, $n$ and $m$ are preserved under the braid move. 

         \item $\ell(w_2)>1$, and $w_2=s_{i+2}s_{i+1}\hdots$, with the dots possibly representing an empty string. Then we have 
        \[w=\hdots\underbrace{\hdots s_{i+2}}_\downarrow\underbrace{s_is_{i+1}}_\uparrow\underbrace{s_is_{i+2}s_{i+1}\hdots}_\downarrow\hdots,\]
        and applying the braid relation we get
        \[v=\hdots\underbrace{\hdots s_{i+2}s_{i+1}s_i}_\downarrow\underbrace{s_{i+1}s_{i+2}}_\uparrow\underbrace{s_{i+1}\hdots}_\downarrow\hdots,\]
        so clearly again the number of elements in an increasing subword is unchanged under the braid move. So again, $n$ and $m$ are preserved under the braid move. 


        \item $\ell(w_2)>1$, and $w_2=s_{i+2}s_i\hdots$, with the dots possibly representing an empty string. Then we have 
        \[w=\hdots\underbrace{\hdots s_{i+2}}_\downarrow\underbrace{s_is_{i+1}}_\uparrow\underbrace{s_is_{i+2}}_\downarrow\underbrace{s_i\hdots}_\uparrow\hdots,\]
        and applying the braid relation we get
        \[v=\hdots\underbrace{\hdots s_{i+2}s_{i+1}s_i}_\downarrow\underbrace{s_{i+1}s_{i+2}s_i\hdots}_\uparrow,\]
        and the number of elements in an increasing subword is unchanged under the braid move. Also, the parity of any end subword is constant. So again, $n$ and $m$ are preserved under the braid move. 








        
        \end{enumerate}
    \end{enumerate}
    This covers all cases, and so the algorithm is constant under braid relations. 
\end{proof}

\end{document}
